
\chapter*{Resumen}
En la actualidad, la interacción virtual entre personas sordas y oyentes se limita casi exclusivamente a la comunicación escrita o a la transcripción unilateral de voz a texto, lo que genera una marcada 
asimetría en la comunicación expresiva durante videollamadas. Este Trabajo Final de Grado desarrolla un prototipo de traducción bidireccional entre la Lengua de Señas Argentina (LSA) y el español oral 
estructurado bajo una arquitectura híbrida de computación en el borde y servicios en la nube con procesamiento en tiempo pseudo-real.

El sistema articula dos canales comunicacionales complementarios. En el sentido señante-oyente, la computadora del usuario captura localmente el flujo de video desde la cámara web, extrae coordenadas 
articulares tridimensionales mediante visión computacional y clasifica un vocabulario de 100 señas mediante una red neuronal ligera (Conv1D + Transformer con \textit{attention pooling}). Con el fin de evitar 
la saturación de memoria y cómputo que provocaría ejecutar modelos de lenguaje generativos en dispositivos convencionales, las glosas clasificadas son enviadas hacia un backend semántico en la nube que 
transforma la secuencia en oraciones fluidas en español rioplatense, las cuales retornan al cliente para su proyección en subtítulos o síntesis de voz. En sentido inverso (oyente-señante), el sistema procesa 
la voz del interlocutor oyente mediante un módulo semántico inverso en la nube que estructura las glosas, enviándolas a un avatar tridimensional en el cliente que ejecuta las animaciones correspondientes o 
recurre al deletreo dactilológico en términos fuera del léxico.

El módulo clasificador gestual local alcanza un 98{,}28\,\% de exactitud top-1 en validación fuera de línea y un 77{,}3\,\% de acierto top-1 (94{,}8\,\% top-3) en cámara real. Por su parte, los modelos de 
lenguaje evaluados (Qwen2.5, Llama 3.2 y SmolLM adaptados mediante LoRA y cuantizados a GGUF Q4\_K\_M) superan el 94\,\% de exactitud y alcanzan latencias de inferencia de entre 200 y 450\,ms. Como mecanismo 
de integración para videollamadas, el sistema se acopla a Google Meet mediante una extensión de navegador conectada al motor local ILSA, asegurando la privacidad del usuario al no transmitir video crudo fuera 
del dispositivo cliente. Finalmente, se deja documentada la propuesta de diseño de una arquitectura en la nube escalable para soportar la demanda concurrente en entornos de producción.

\chapter*{Palabras clave:} Lengua de Señas Argentina, reconocimiento de gestos, arquitectura híbrida cliente-nube, traducción semántica bidireccional, avatar 3D, accesibilidad.

% TODO: Chequear donde se usan por primera vez, si conviene agregarlas al resumen
% LSA
% LLM
% ASL
% SML